\chapter{Lint, simulation and formal}
\label{ch:flows}

After \cmd{bag generate}, three commands check the generated design with EDA tools:
\cmd{bag lint}, \cmd{bag sim} and \cmd{bag formal}. Each takes the output directory as
its only argument and reads \file{bag_outputs.json} to find the files, the top module
and the settings (section~\ref{sec:manifest}). \cmd{bag flow} runs the whole chain from
the spec. The tools must be installed at the versions of section~\ref{sec:eda}; the
tools image has all of them.

\section{How a report reads}
\label{sec:report}

Every flow prints a \term{report}. A report has one line per \term{finding}, then a
summary line, then any notes. A finding line reads:

\begin{center}
\code{TOOL: FILE:LINE: SEVERITY: MESSAGE [RULE]}
\end{center}

\code{TOOL} is the tool that reported it. \code{FILE} is relative to the output
directory; a finding that has no place leaves out \code{FILE:LINE:}. \code{SEVERITY} is
\code{error} or \code{warning}. \code{RULE} is the tool's own name for the finding, or
one of BAG's rule IDs below. A finding that a waiver covers ends with
\code{(waived)}.

The summary line names the flow and its tools, and counts the findings:
\code{N unwaived finding(s), M waived}. A flow passes when it has no unwaived finding,
error or warning. Then the command exits with status 0. Otherwise it exits with status
1. A note (\code{note: ...}) is information for you, for example that a tool is not the
pinned version. A note never makes a flow fail.

Two rule IDs can come from any tool:

\begin{description}
\item[\ruleid{TOOL-MISSING}] The tool's program is not on your \code{PATH}. Install it,
  or run BAG in the tools image.
\item[\ruleid{TOOL-FAILED}] The tool ended with a non-zero exit status but reported no
  finding that BAG could read. The message gives the exit status and the tool's last
  line of output.
\end{description}

\section{Lint}
\label{sec:lint}

\cmd{bag lint} checks the generated RTL with static tools. It runs the tools of Project >
\setting{Lint tools}, in that order, on every file of the file list, with the top of the
manifest as the top module. The default is all three tools:

\begin{description}
\item[slang] \code{slang --quiet -Weverything}, plus the options of BAG's
  \file{slang.flags} (below). It elaborates the whole design.
\item[Verilator] \code{verilator --lint-only -Wall}. It elaborates the whole design.
\item[Verible] \code{verible-verilog-lint} with Verible's default rules. It checks style,
  one file at a time, and skips the vendored files.
\end{description}

The lint passes when no tool reports a finding that is not waived. A warning counts as
much as an error. Verilator reports an elaboration-time \code{$error} or \code{$fatal}
in the RTL (\code{USERERROR}, \code{USERFATAL}) as a warning; BAG judges it as an error,
because it is the design's own check of an illegal parameter.

\lstinputlisting[style=console]{examples/console/eda/flow-lint.txt}

\subsection{Third-party RTL}

The files under \file{rtl/vendor/} are \term{third-party} RTL: the PULP files and the
files BAG derives from them. slang and Verilator still read them, so that every
instantiation of a vendored module is checked. But BAG does not judge the warnings in
those files: it counts them in one line, \code{N third-party warning(s) in rtl/vendor/ not
judged}, and the lint still passes. An error in a vendored file still fails the lint.

A Verilator \code{UNOPTFLAT} warning (combinational logic that loops) is reported at one
signal of the loop. When the rest of the loop is in vendored files or in BAG's
in-house wrappers around them, BAG moves the warning to the first vendored file on the
loop and counts it as third-party. A loop that runs through any other generated signal
stays where it is and is judged: that is where a generator error could wire a READY
back into a VALID.

\subsection{Built-in configuration and waivers}

BAG ships its lint configuration. Each option and each waiver carries its reason.
A \term{waiver} accepts one kind of finding at the places it names, so that the finding
does not fail the lint.

The options that \file{slang.flags} adds to \code{-Weverything}:

\begin{description}
\item[\code{-Wno-unused-def}] The file list holds whole libraries, so some modules in it
  are not used by this design.
\item[\code{-Wno-case-redundant-default}] Every \code{case} lists its enum values and
  keeps a \code{default} for X and reset safety.
\item[\code{-Wno-useless-cast}] A width cast in parameterised RTL is an identity for some
  parameter values.
\item[\code{-Wno-unused-but-set-variable}] The \code{unused} sink idiom; Verilator's
  \code{UNUSED} warnings check unused signals instead.
\end{description}

Verible runs its default rules with no change. The built-in waivers are:

\begin{center}
\begin{tabular}{lll}
\toprule
Tool & Rule & Files \\
\midrule
slang & \code{unused-package-parameter} & \file{rtl/pkg/*.sv} \\
Verilator & \code{UNUSEDPARAM} & \file{rtl/pkg/*.sv} \\
\bottomrule
\end{tabular}
\end{center}

Both waivers have the same reason. The parameter package exports every fabric table and
the whole memory map for the SoC that instantiates the wrapper; the wrapper itself reads
only the tables that its templates need. So most designs show some waived findings in
the package. There is no Verible waiver.

A built-in waiver whose files are in the design but that waives nothing is itself an
error, \ruleid{WAIVER-UNUSED}, so that stale waivers cannot pile up. BAG reports it only when the tool ran to the end (no
error stopped it early) and is the pinned version that the waivers were written for.
With another version it is a note (section~\ref{sec:versions}).

You cannot add waivers of your own from the command line.

\section{Compile}
\label{sec:compile}

The compile step elaborates the design with Verilator as a simulator would: \code{verilator
--cc} with the file list and the top of the manifest. Lint warnings are switched off here,
because they are the lint step's job; only errors count. The model goes to
\file{obj/<top>/} in the output directory.

There is no command of its own for this step: \cmd{bag flow} runs it after the lint. When
the design has generated tests, it only elaborates (\code{verilator-elaborate} in the
report), because \cmd{bag sim} builds the simulation model itself. Without generated
tests, it also builds the C++ model (\code{verilator-build}), so that the design is
compiled at least once.

\section{Simulation}
\label{sec:sim}

\cmd{bag sim} runs the generated cocotb tests (section~\ref{sec:tests}) with Verilator. It
needs the \code{sim} extra of BAG, which installs cocotb and cocotbext-axi
(chapter~\ref{ch:install}). It builds the model in \file{sim/} of the output directory,
runs every test of \file{tests/<project>_tests.py}, and reads the results that cocotb
writes (\file{sim/results.xml}).

\cmd{bag sim} stops with an error, before it runs anything, when:

\begin{itemize}
\item the output directory has no generated tests (Project > \setting{Generate cocotb unit
  tests} is N, or the design has a user-supplied fabric);
\item Project > \setting{Simulator} is not \code{verilator};
\item cocotb is not installed.
\end{itemize}

The run also fails when the output directory was generated by another version of BAG:
generate it again with the version you run.

\lstinputlisting[style=console]{examples/console/eda/flow-sim.txt}

\subsection{What the tests check}

The tests drive the top through bus models: one model per master port and one per slave
port. A \term{reference model} in BAG computes, from the IR, what every access must do:
which slave it reaches and at which address, or which error it gets. Each test compares
what the design does with the reference model. Every test also counts what it checked,
and fails if it checked nothing.

BAG writes only the tests that can check something on your design. The full list, in the
order they run:

\begin{description}
\item[\code{every_master_reaches_what_the_model_says}] Every master writes and reads the
  first, middle and last words of every region, the words around it, and both edges of
  every alias window, with each setting of the alias pins and, where there are secure
  regions, both security states. Each access reaches the slave and address that the
  model names, or gets the error it names. Every design has this test.
\item[\code{arbitration_follows_the_fabric_policy}] AHB-Lite fabrics with contention: the
  arbiter grants in the order of the fabric's policy.
\item[\code{concurrent_traffic_keeps_every_masters_data}] All masters at once, with
  several transactions in flight: every write lands, and every read returns that
  master's own data.
\item[\code{traffic_survives_back_pressure_on_every_channel}] The same traffic, single
  beats and then bursts, with random stalls on every channel: at the slaves, then at the
  masters, then at both.
\item[\code{reset_during_traffic_leaves_nothing_behind}] Every reset is asserted while a
  request to each slave is in its handshake. While reset is held, nothing asks a model
  to act. After it, nothing left over appears, and fresh traffic returns the right data.
\item[\code{traffic_crosses_clock_domains_at_every_ratio}] Designs with clock-domain
  crossings: traffic at several clock ratios (equal periods out of phase, 1:3, 3:1, 10:7
  and 7:10).
\item[\code{shared_buses_block_behind_a_stalled_slave_and_crossbars_do_not}] With one
  slave stalled, a second master's access waits on a shared bus and completes on a
  crossbar.
\item[\code{slave_decerr_returns_as_the_model_says}] An AXI slave answers DECERR: each
  master gets the response that its path gives it, and the same word works afterwards.
\item[\code{exclusive_accesses_carry_lock_and_exokay_unchanged}] AXI4 exclusive accesses:
  AxLOCK reaches the slave and EXOKAY reaches the master. BAG does not implement
  exclusive accesses; it must not disturb them.
\item[\code{axi_arbitration_rounds_starvation_and_lock}] AXI fabrics with several masters
  on one slave: the order of grants per round, starvation, and that a grant does not
  change before its handshake.
\item[\code{axi_shared_bus_midpoint_buffers_reads}] An AXI shared bus with a mid-point
  FIFO or register slices: a master can issue exactly as many reads as the mid-point
  holds.
\item[\code{axi4_bursts_reach_slaves_and_error_paths}] AXI4 INCR bursts reach their slave,
  and a burst to an unmapped or denied address gets the error on every beat.
\item[\code{axi4_crossbars_admit_max_mst_trans_per_id}] AXI4 crossbars admit
  Fabrics > \col{Max mst trans} reads on one ID.
\item[\code{narrow_apb_behind_axi2apb_keeps_lanes_strobes_and_errors}] APB slaves narrower
  than 32 bits behind \code{bag_axi2apb}: addresses, byte lanes, strobes, and an error on
  any beat.
\item[\code{unaligned_axi_writes_store_only_their_bytes}] A write that starts 1, 2 or 3
  bytes into a word stores exactly its bytes.
\item[\code{an_error_on_any_beat_of_a_split_access_reaches_the_master}] A word split into
  several beats for a narrow slave: an error on any beat reaches the master.
\item[\code{axi4_burst_shapes_match_the_model}] Every AXI4 burst shape that the model
  knows, written and read back; where a converter on the path refuses the burst, the
  master gets SLVERR and nothing is stored.
\item[\code{single_bytes_reach_both_ends_of_every_region}] AXI4 and AHB-Lite masters:
  single bytes at the first and last byte of every region, and just outside it.
\item[\code{apb_byte_writes_carry_pstrb}] A byte write to an APB4 slave carries the right
  PSTRB, and the other bytes keep their value.
\item[\code{axi_ports_take_every_handshake_ordering}] On every AXI channel that a model
  answers: READY before VALID, after VALID, and in the same cycle. A design whose VALID
  waits for READY times out.
\end{description}

A test fails when the design does something else, or when it runs out of time; each test
has a time limit, so a design that deadlocks fails instead of hanging. The rule IDs of
\cmd{bag sim}, each an error:

\begin{description}
\item[\ruleid{TEST-FAILED}] A test failed. The message names it.
\item[\ruleid{TEST-SKIPPED}] A test was skipped. A generated test must run.
\item[\ruleid{TEST-SET}] The tests that ran are not the tests that were expected.
\item[\ruleid{SIM-EXIT}] Every test passed, but the simulator then exited with a non-zero
  status.
\end{description}

\subsection{The seed}

The tests make random choices, for example the stall patterns. They take them from one
fixed \term{seed}, \code{20261031}, so a run gives the same result every time and a
failure can be repeated. Every finding of \cmd{bag sim} names the seed. To try another seed, set the
environment variable \code{BAG_COCOTB_SEED} to a number before you run \cmd{bag sim}.
One seed is a limit of the tests (chapter~\ref{ch:limits}).

\subsection{Running some of the tests}

\opt{--filter} takes a regular expression and runs only the tests whose name it finds.
The name it searches is \code{<project>_tests.<test>}, for example
\code{lite_ss_tests.every_master_reaches_what_the_model_says}. An expression that finds
no test is an error.

\subsection{The READY shim}

\opt{--ready-shim} runs the test \code{traffic_survives_back_pressure_on_every_channel}
through the READY shim of \file{tests/} instead of the top (section~\ref{sec:tests}). The
shim drives each AXI READY into the design from that channel's VALID. If the design
waits for READY before it raises VALID, the loop does not settle and the test fails. The
build goes to \file{sim_ready_shim/}. The option cannot be used with \opt{--filter}, and
it is an error when the design has no shim.

\lstinputlisting[style=console]{examples/console/eda/flow-ready-shim.txt}

\subsection{The coverage report}

\cmd{bag sim} writes a \term{functional coverage} report, \file{sim/coverage.md}, and the
counts it is made from, \file{sim/coverage.json} (in \file{sim_ready_shim/} with
\opt{--ready-shim}). A \term{bin} is one thing the spec says
the design does. The report has four kinds of bins:

\begin{description}
\item[path] each master reaches each slave, by read and by write;
\item[response] each response kind a master gets from each slave: OKAY, a slave error, or
  a miss;
\item[secure] each master's accesses to each secure-only region, secure and non-secure,
  allowed or refused;
\item[burst] each AXI4 master's burst shapes, and whether a converter on the path
  refuses them.
\end{description}

The report of \cmd{bag sim} ends with one note per kind, \code{coverage: KIND COVERED/BINS},
and a note with the path of the coverage report. A test records a bin only after it has checked that
access. So a run with \opt{--filter}, or a test that fails, leaves bins uncovered.

The coverage report is a report only. It has no closure target, and it never makes
\cmd{bag sim} fail. Its bins come from the same addresses and burst shapes that the
tests use, so a full passing run always reads 100\%: the report shows what a run
exercised, not what the tests never try (chapter~\ref{ch:limits}).

The transcript in section~\ref{sec:sim} ends with the coverage report of the
\code{lite_ss} run.

\section{Formal}
\label{sec:formal}

\cmd{bag formal} runs the SymbiYosys jobs that \cmd{bag generate} wrote into
\file{formal/} (section~\ref{sec:jobs}), one job after the other, and prints one report
for all of them. A job runs in \file{formal/}; each of its tasks gets a work folder there,
\file{formal/<job>_<task>/}, which the next run replaces.

\subsection{Jobs, tasks and modes}

A \term{job} is one \file{.sby} file. It has one or more \term{tasks}; SymbiYosys names
each task \code{<job>_<task>}, for example \code{lite_ss_reset_prove}. Each task runs in
one of three \term{modes}, to a \term{depth} counted in \term{solver steps}:

\begin{description}
\item[\code{bmc}] A \term{bounded} check. The solver tries every input sequence of up to
  \emph{depth} steps from the start, and looks for one that makes an assertion fail. A
  pass means that no assertion fails within that many steps. It says nothing about a
  later step. The report adds a note for each bmc task:
  \code{<task>: bounded check to depth N (bmc), not a proof}.
\item[\code{prove}] A proof by induction (k-induction with \emph{depth} as k). A pass
  means that the assertions hold in every state the design can reach, however long it
  runs.
\item[\code{cover}] The solver looks for an input sequence of up to \emph{depth} steps
  that reaches each \code{cover} statement. A cover shows that the state an assertion
  checks really happens. A cover that is not reached means the assertion may hold only
  because its condition never occurs (it holds \term{vacuously}). That is a finding.
\end{description}

The clocks of a design with several clock domains are free inputs, so every clock ratio
and phase is checked. Then each clock edge takes about two solver steps. A design or
unit with one clock domain runs on one clock, and each solver step is one edge.

\subsection{The jobs}

\begin{description}
\item[\file{<project>_reset}] Quiet in reset: while a domain's reset is held, the design
  asks nothing of the masters and slaves in that domain. No request VALID, PSEL or
  HTRANS goes to a slave, no response VALID to an AXI master, and HREADY is high towards
  an AHB-Lite master. Tasks: \code{prove} (depth 6) and \code{cover} (depth 40). This
  is a proof: it holds in every reachable state.
\item[\file{<project>_protocol}] Port protocol: each port of the top follows the rules of
  its protocol. On a port where the design is the slave, the checker asserts the rules
  for what the design drives and assumes them for what the environment drives; on a
  port where the design is the master, the other way round. AXI4-Lite and APB ports use
  the ZipCPU checkers, AXI4 ports \file{bag_faxi4.sv} and AHB-Lite ports
  \file{bag_fahb.sv}. The checks that a response comes within a bounded time are off,
  because a port waits as long as the slaves beyond it stall. Tasks: \code{bmc} (depth
  16), and a \code{cover_<port>} task for each port, at a depth that fits the path from
  that port to a peer (at least 16). A port whose covers its fabric's route job already
  holds has no cover task here. This is a bounded check: a cover reached after step
  16 shows that the check can fire there, not that the bmc checked that far.
\item[\file{<project>_decode}] Address decode: for every fabric that BAG generates, the
  address decoder with the fabric's tables from the parameter package is compared with a
  reference that BAG writes from the IR: the aliases, the regions with their access and
  secure-only flags, the \sheet{Access_Rules} sheet, the miss policy and the fabric's address
  window. The reference does not read the package, so a wrong table entry shows as a
  difference. The inputs are free: the master, the address, read or write, the secure bit
  and the alias-enable pins. Task:
  \code{prove} at depth 1: the check is combinational, so one step covers every input,
  and the result holds for every address, master, direction and secure bit.
\item[\file{<project>_route_<unit>}] One job per fabric and per bridge, checked on its own
  with free inputs. For a fabric: every request at a slave port is one that a master made
  and that the reference sends to that port; a miss or a denied access reaches no slave
  port; and a master gets a response only for a request that a slave took or that the
  fabric answers itself. For a fixed-priority shared bus or AHB-Lite bus: of two
  requests that start at the same edge, the one with the lower Priority value goes first.
  For a bridge: every downstream request lies in the window of the upstream address.
  Tasks: \code{bmc} and \code{cover}, both to a depth that depends on the unit (its clock
  domains, width conversions, FIFOs and register slices). This is a bounded check.
\end{description}

\subsection{What the jobs assume}

A formal check holds under the \term{environment} that it assumes: the rules that the
masters, slaves, clocks and resets outside the design follow. The generated jobs assume
the following.

\begin{itemize}
\item \textbf{Resets.} Every reset is asserted at the start. A reset asserts only together
  with the rest of its reset group, and the resets of a group release in any order. A
  reset is released only after it has taken effect. This is the reset sequence of
  \file{doc/reset.md}. In a route job, each domain starts in reset and is then released
  for good.
\item \textbf{Quiet in reset.} While their domain is in reset, the masters and slaves
  drive no VALID (AXI) and no PSEL (APB), and drive HTRANS IDLE and HREADYOUT high
  (AHB-Lite). They stay quiet until the clock edge after the reset is released, as AMBA
  requires of an AXI master, and until every domain of the reset group has left reset.
\item \textbf{Handshakes.} What the environment drives into a port follows the rules of
  that port's protocol. The protocol checkers assume it.
\item \textbf{Alias pins.} An alias-enable pin (\sheet{Address_Decode} > \col{Alias enable})
  changes only while the design is in reset.
\end{itemize}

If your SoC breaks one of these rules, the results do not apply to it. Alias pins that
must not change outside reset are a limit of the design (chapter~\ref{ch:limits}).

\lstinputlisting[style=console]{examples/console/eda/flow-formal.txt}

\subsection{What the findings mean}

Every formal finding is an error.

\begin{description}
\item[\ruleid{FORMAL-ASSERT}] An assertion fails. The finding names the assertion, its
  file and line, and the counterexample trace: a waveform file in the task's work folder
  that shows the input sequence that breaks it.
\item[\ruleid{FORMAL-INDUCTION}] In a \code{prove} task, an assertion holds to the depth
  but the induction does not close: the solver found a state, maybe not reachable, from
  which the assertion fails. The finding names the trace. An assertion that fails
  outright is reported as \ruleid{FORMAL-ASSERT} only.
\item[\ruleid{FORMAL-COVER}] A cover is not reached within the depth, so the assertion
  that it guards may hold only vacuously.
\item[\ruleid{FORMAL-TIMEOUT}] The job gave no verdict within the time that
  \opt{--timeout} allows. The finding names the last solver step reached, if the tool
  printed one. For a cover
  task, a note lists the covers reached before the limit.
\item[\ruleid{FORMAL-ERROR}] The tool stopped on an error: a parse or elaboration error,
  or assumptions that allow no input sequence at all. In that last case nothing was
  checked, and the message says so.
\item[\ruleid{FORMAL-UNKNOWN}] The task ended with a status that is not a verdict BAG
  knows (for example UNKNOWN), or failed without naming a property.
\item[\ruleid{TOOL-FAILED}] SymbiYosys exited with an error and printed nothing that BAG
  could read, or a task of the job printed no verdict at all.
\item[\ruleid{TOOL-MISSING}] \code{sby} is not on your \code{PATH}.
\end{description}

\subsection{Time limits}

Formal checks can take long. On a large design, a job can take many minutes, and the
cover tasks the longest. Without \opt{--timeout}, \cmd{bag formal} waits for each job as
long as it takes. \opt{--timeout} sets the seconds allowed for each job (at least 1),
with all its tasks. When the time is up, BAG stops the job, reports \ruleid{FORMAL-TIMEOUT}, and goes on with the next job.

\cmd{bag formal} stops with an error, before it runs anything, when the output directory
has no formal jobs (Project > \setting{Generate formal (SymbiYosys)} is N, or the design
has a user-supplied fabric), or when a job file that the manifest lists is missing.

\section{The whole flow}
\label{sec:flow}

\cmd{bag flow} takes a spec, not an output directory, and runs every step in turn:

\begin{enumerate}
\item check and generate, as \cmd{bag generate} (with its \opt{-o} and \opt{--style});
\item lint, as \cmd{bag lint};
\item compile (section~\ref{sec:compile});
\item the generated tests, as \cmd{bag sim}, if the design has them; otherwise a line
  says why there are none;
\item with \opt{--formal}, the formal jobs, as \cmd{bag formal}; again a line says so when
  there are none.
\end{enumerate}

Each step prints its report. The flow stops at the first step that fails, with exit
status 1. Formal is left out unless you give \opt{--formal}, because the proofs of a
large design can take many minutes; you can run \cmd{bag formal} on the output directory
later instead. \opt{--formal-timeout} is the \opt{--timeout} of the formal step, and needs
\opt{--formal}.

\lstinputlisting[style=console]{examples/console/eda/flow-flow.txt}

\section{Random configurations}
\label{sec:regress}

\cmd{bag regress} tests BAG itself on configurations nobody wrote by hand. For each seed
it draws a configuration (fabrics of every protocol and topology, bridges, adapters,
address maps, aliases, hierarchy and port options), writes its workbook, and runs it
through every step. A seed draws the same configuration on every machine, so a result can
be repeated with \cmd{bag regress --seed S}.

Each seed ends with one status:

\begin{description}
\item[pass] Every step passed and every formal task ran to a verdict.
\item[fail] A step reported a finding or stopped with an error. \code{step} in
  \file{result.json} names it.
\item[timeout] Formal work was cut off by \opt{--formal-budget} or not reached. Formal runs
  task by task: the decode proof first, then the reset proof, the route and protocol
  bounded checks, then the covers. After the decode job each task may use a quarter of the
  budget left (at least 30 seconds). A timeout is reported, not a failure.
\item[generator] The drawn configuration fails \cmd{bag check}: a bug in the generator,
  not in the design.
\end{description}

\file{result.json} also records the findings, the dimensions the configuration covers
(protocol pairs, width ratios, clock crossings and so on), the budget and the version of
BAG, and the commands that repeat it: \cmd{bag regress --seed S} and \cmd{bag flow} on its
IR. A seed's configuration depends on the version of BAG, so repeat a seed with the same
version.

\section{Tool versions}
\label{sec:versions}

BAG is tested with one version of each tool, the version pinned in \file{tools.lock}
(section~\ref{sec:eda}). The built-in waivers and the way BAG reads each tool's output
are written for those versions. BAG asks each tool for its version, and says so in the
report when it is not the pinned one:

\begin{itemize}
\item \cmd{bag formal} adds a note when \code{sby} is not v0.69:
  \code{sby <version> is not the pinned v0.69}.
\item \cmd{bag lint} adds a note when a built-in waiver waived nothing and the tool is
  not the pinned version: the note names the tool, its version, the pinned version and
  the waivers. With the pinned version the same waiver is the error
  \ruleid{WAIVER-UNUSED}.
\end{itemize}

The simulation and compile steps do not check the version of Verilator. Another version
of a tool can also report findings that the pinned version does not. If you see a
finding that CI does not, compare the versions first.
